\section{Zusammenfassung}

Ziel der vorliegenden Arbeit war es, die Balanced Scorecard als strategisches Managementinstrument darzustellen und auf die Brau Union Österreich AG (BUÖ) anzuwenden. Das Unternehmen nimmt mit einem Marktanteil von über 50 \% eine führende Position im österreichischen Biermarkt ein, steht jedoch gleichzeitig vor Herausforderungen wie einem zunehmend gesättigten Markt und veränderten Konsumgewohnheiten. Daher ist es erforderlich, die Unternehmensstrategie in konkrete Ziele, Kennzahlen und Maßnahmen zu übersetzen.

Die Balanced Scorecard bietet hierfür ein geeignetes Instrument, da sie strategische Ziele nicht nur aus finanzieller Sicht betrachtet, sondern auch Kunden, interne Prozesse sowie Lernen und Entwicklung einbezieht. Durch diese mehrdimensionale Perspektive können strategische Vorgaben systematisch in operative Maßnahmen überführt und deren Umsetzung besser gesteuert werden.

Auf dieser Grundlage wurde eine Balanced Scorecard für die BUÖ entwickelt. In der Finanzperspektive stehen insbesondere die nachhaltige Steigerung des Unternehmenswertes, ein Return on Investment oberhalb der Kapitalkosten sowie die kontinuierliche Erhöhung des Free Cash-flows im Mittelpunkt. Diese Ziele bilden den finanziellen Rahmen für die langfristige Unternehmensentwicklung.

In der Kundenperspektive liegt der Fokus auf der Sicherung und dem Ausbau der Marktführerschaft, der Maximierung der Kundenzufriedenheit im Lebensmittelhandel und in der Gastronomie sowie der Erschließung neuer Kundensegmente. Dadurch sollen sowohl bestehende Marktanteile gesichert als auch neue Wachstumspotenziale erschlossen werden.

Die interne Prozessperspektive konzentriert sich auf effiziente Produktions-, Logistik- und Vertriebsprozesse sowie auf die Sicherstellung hoher Qualitätsstandards. Diese Prozesse bilden die Grundlage für eine zuverlässige Belieferung der Kunden und tragen wesentlich zur Wettbewerbsfähigkeit des Unternehmens bei.

Ergänzend dazu zeigt die Lern- und Entwicklungsperspektive die Bedeutung der Mitarbeiter und der Innovationsfähigkeit für den langfristigen Unternehmenserfolg. Die Förderung unternehmerisch denkender Mitarbeiter, eine stärkere bereichsübergreifende Zusammenarbeit sowie Innovationen im Produktportfolio schaffen wichtige Voraussetzungen für eine nachhaltige Weiterentwicklung des Unternehmens.

Insgesamt zeigt die entwickelte Balanced Scorecard, wie die strategischen Ziele der BUÖ in ein strukturiertes System aus Kennzahlen und Maßnahmen überführt werden können. Durch die Verknüpfung der verschiedenen Perspektiven wird eine ganzheitliche Steuerung des Unternehmens ermöglicht und die langfristige Wettbewerbsfähigkeit unterstützt.